\begin{formbox}
\textbf{Name:} Jiasheng Xu \hfill \textbf{Student ID:} 2251640\\
\textbf{Role:} Testing and Validation (Role 5).

\textbf{1. My specific contributions to this project:}
\begin{enumerate}[leftmargin=2em,itemsep=0.25em,topsep=0.15em]
    \item I added the independent test suite in
    \path{project2/test/test_validation.py} for the team's variable-diffusion
    finite-difference and P1 finite-element solvers. I tested asymmetric,
    nonzero Dirichlet data and node ordering, and ran course-bank cases
    P1, P2, P4 and P5 through both assemblers on four meshes, for 32 solves.
    The suite passed nine tests with the bank enabled. The original
    \path{run_all.py} entry point passed all seven of its tests.

    \item I used the same P1-reconstructed \(L^2\) norm for both methods.
    Fitting the finest three of four meshes gave orders of \(2.000\) for
    FDM and \(1.999\) for FEM. At \(N=32\), tightening the conjugate gradient
    tolerance to \(10^{-12}\) brought both solutions within
    \(1.1\times10^{-14}\) of sparse LU in the nodal maximum norm. The
    discretisation errors were about \(8\times10^{-4}\).

    \item I recorded the commands, measured errors and reproducible failures
    in \path{project2/test/VALIDATION_RUN_2026-11-06.md}
    (commit \texttt{08bf830}). A closed-form spectral check showed that the
    condition-number routine returned \(89.045\) for the uniform-material
    case at \(N=15\), where the analytic value is \(90.523\) to three
    decimal places. The team's results section uses the corrected estimates.
\end{enumerate}

\textbf{2. The hardest conceptual obstacle I encountered:} The main challenge
was finding a diagnostic that could detect incorrect interface resistance.
At \(c=10\) and \(N=63\), arithmetic face averaging produced a reconstructed
two-sided flux jump of about \(4\times10^{-13}\). The maximum shared-face flux error against the
analytic flux was \(0.662\). I included both measurements in the validation
so the small jump could be interpreted alongside the physical flux error.

\textbf{3. One piece of AI-generated output my team used, and what I changed:}
I used Codex to draft and run the independent validation scripts. I asked
for the course-bank cases to be solved with both assemblers, for the linear
residual to be recomputed, and for the interface flux to be checked against
the exact solution. I also requested the Project 1 test-directory layout
so that the revised checks could be rerun with \texttt{unittest} and traced
to the measured logs.

\textbf{4. One thing I still do not fully understand:} I would like to
derive a bound for the maximum interface-flux error near the points where
the material interface meets the Dirichlet boundary, and understand how that
bound depends on mesh size and material contrast.

\vfill
\textbf{Signature:} \underline{\hspace{4.5cm}} \hfill
\textbf{Date:} \underline{\hspace{2.5cm}}
\end{formbox}
